% !TeX TS-program = xelatex
\documentclass{resume-photo}

\usepackage{xeCJK}
\setCJKmainfont{Noto Serif CJK SC}
\usepackage{fontspec}
\usepackage{geometry}
\geometry{top=0.55cm, bottom=0.45cm, left=0.85cm, right=0.85cm, includefoot}
\usepackage{enumitem}
\setlist[itemize]{itemsep=0pt, topsep=1pt, parsep=0pt, partopsep=0pt}
\linespread{1.05}

\ResumeName{赵伟}

\begin{document}

\ResumeContacts{
  135-XXXX-9963,%
  \ResumeUrl{mailto:zhaowei@bupt.edu.cn}{zhaowei@bupt.edu.cn},%
  \textnormal{北京邮电大学 | 计算机科学与技术 · 本科 | 20XX-XX}%
}

\ResumeTitle

% ====================== 简介 ======================
\noindent {\small 北京邮电大学计算机科学与技术本科毕业，5 年 Java 后端开发经验，专注高并发、高可用分布式系统设计与落地。先后在互联网与金融科技领域参与核心交易系统建设，具备亿级流量服务治理经验，熟悉 Spring Cloud 微服务全技术栈与中间件调优。}

% ====================== 教育 ======================
\section{教育背景}
\ResumeItem{北京邮电大学}[\textnormal{计算机学院 | 计算机科学与技术 | 本科 | 专业排名前 8\%}][2015.09—2019.06]

% ====================== 技能 ======================
\section{专业技能}
\begin{itemize}
  \item \textbf{语言/框架}：Java 8/11/17、Kotlin；Spring Boot / Spring Cloud / Dubbo；熟悉响应式编程（WebFlux / Reactor）。
  \item \textbf{存储/缓存}：MySQL（InnoDB 调优、分库分表 ShardingSphere）；Redis 集群；Elasticsearch；HBase。
  \item \textbf{消息/流}：Kafka（百亿级消息吞吐）；RocketMQ；熟悉 Flink 流处理与实时数仓。
  \item \textbf{基础设施}：Kubernetes / Docker；Prometheus + Grafana 监控；Skywalking 链路追踪；熟悉 AWS / 腾讯云。
\end{itemize}

% ====================== 工作 ======================
\section{工作经历}

\ResumeItem{蚂蚁集团｜支付宝核心交易平台}[\textnormal{高级 Java 工程师}][2022.03—至今]
\begin{itemize}
  \item 参与收单交易核心链路重构（Dubbo 3.x 升级），服务调用 P99 延迟从 120ms 降至 45ms，双十一峰值 TPS 支撑 15 万+。
  \item 设计幂等框架（分布式锁 + 唯一流水号 + 状态机），解决网络重试导致的重复扣款问题，资损事件归零。
  \item 主导 MySQL 分库分表改造（ShardingSphere），将订单表从单库 8 亿行拆分至 16 库 256 表，查询耗时降低 85\%。
  \item 构建实时对账流水线（Kafka + Flink），账单核对时效从 T+1 批处理缩短至准实时 5 分钟，差错发现率提升 40\%。
\end{itemize}

\ResumeItem{滴滴出行｜出行服务后端团队}[\textnormal{Java 工程师}][2019.07—2022.02]
\begin{itemize}
  \item 负责行程匹配服务后端开发，基于 Redis GEO + 自研调度算法实现司机就近匹配，响应时间 < 30ms，高峰 QPS 50 万+。
  \item 参与动态定价模块开发，引入 Flink 实时计算供需比，定价策略迭代周期从周级缩短至分钟级。
  \item 推进核心服务容器化（K8s），节省 30\% 机器资源；基于 Skywalking 建立全链路监控，故障定位时间从 2h 降至 10min。
\end{itemize}

% ====================== 项目 ======================
\section{项目经历}

\ResumeItem{开源中间件：easy-distributed-lock}[][2021.06—至今]
\begin{itemize}
  \item 基于 Redis / Zookeeper 实现支持可重入、公平锁、读写锁的分布式锁框架，兼容 Spring Boot Starter 自动装配。
  \item GitHub 获得 900+ stars，Maven Central 月下载量 3000+；被 10+ 家企业用于生产环境。
\end{itemize}

% ====================== 荣誉 ======================
\section{个人荣誉}
\begin{itemize}
  \item 蚂蚁集团 2023 年度"卓越工程师"（支付宝核心交易方向 Top 5\%）
  \item 北京邮电大学 2019 年度优秀毕业生；ACM-ICPC 北京区域赛银奖
  \item InfoQ 认证技术作者；CSDN Java 领域博主（粉丝 2 万+）；掘金年度优秀作者 2022
\end{itemize}

\end{document}
